\begin{document}

%----------HEADING----------
\begin{center}
    {\Huge\bfseries Amani KRID} \\ \vspace{5pt}
    {\large\color{dark-grey} Ingénieure logiciel \& Data} \\ \vspace{7pt}
    \small
    \texttt{kridamani.pro@gmail.com} \hspace{3pt}$|$\hspace{3pt}
    07 81 25 95 21 \hspace{3pt}$|$\hspace{3pt}
    Nice, France \\ \vspace{3pt}
    \href{https://github.com/KRIDAmani}{github.com/KRIDAmani} \hspace{3pt}$|$\hspace{3pt}
    \href{https://www.linkedin.com/in/amani-krid-63aa3723a}{linkedin.com/KRIDAmani}\hspace{3pt}$|$\hspace{3pt}
    \href{https://kridamani.github.io/Personal-Portfolio/}{AmaniKRID/Personal-Portfolio} \hspace{3pt}
    
\end{center}

%-----------PROFIL-----------
\section{PROFIL}
Élève-ingénieure en dernière année à l'École Centrale de Lyon (double diplôme ENSI, \textbf{BAC+6}), je recherche un poste en Data Engineering ou AI Engineering en France à l'issue de mon stage de fin d'études. Basée à Nice, je suis mobile sur tout le territoire.

%-----------EXPERIENCE-----------
\section{EXPÉRIENCES}
\resumeSubHeadingListStart
    \resumeSubheading
      {Amadeus SAS}{Avril 2026 -- Septembre 2026}
      {Aide à la décision par IA pour la reprise nouvelle génération après perturbation aérienne}{Sophia Antipolis, France}
      \resumeItemListStart
        \resumeItem{Conception d'un pipeline de machine learning de bout en bout : classement des solutions (LambdaRank), estimation de coût calibrée (LightGBM) et optimisation combinatoire par programmation linéaire en nombres entiers, remplaçant le moteur heuristique de re-accommodation des passagers utilisé en production par de grandes compagnies aériennes.}
       \resumeItem{Collaboration en méthodologie agile (SAFe) au sein d'une équipe internationale (Nice, Bogota, Sydney) entre l'équipe innovation et l'équipe produit historique, avec présentations régulières des résultats aux parties prenantes.}
        \resumeItem{\tools{Python, PEP8, PuLP/CBC, OR-Tools, HiGHS, scikit-learn, Splunk, SAFe/Scrum, Jira, Confluence, Streamlit, Robot, Azure}}
      \resumeItemListEnd
    \resumeSubheading
      {Laboratoire LIRIS}{Mai 2025 -- Juillet 2025}
      {Mobilité urbaine : dispatch prédictif de taxis}{Lyon, France}
      \resumeItemListStart
        \resumeItem{Simulateur de dispatch de taxis (découpage dynamique de zones, relocalisation proactive, prédiction ARIMA) et adaptation d'une méthode de référence en apprentissage par renforcement ; \textbf{taux de rejet réduit de 28\% à 12\%} face au dispatch aléatoire, attente moyenne \textbf{2,71 min}, sur données réelles (Porto, New York, Lyon).}
        \resumeItem{\tools{Python, NetworkX, OSMnx, ARIMA, KMeans, pandas, GitLab}}
      \resumeItemListEnd
    \resumeSubheading
      {Primatec Engineering -- BMW}{Juillet 2024 -- Août 2024}
      {Transformation YAML $\rightarrow$ JSON pour systèmes embarqués}{Sfax, Tunisie}
      \resumeItemListStart
        \resumeItem{Outil de conversion et de recherche en Rust pour données de calculateurs embarqués : benchmark de 6 stratégies ramenant le temps de recherche de \textbf{4,4\,s à 0,11\,s (\textasciitilde40$\times$)} via générateurs, traitement par lots et recherche plein texte (FTS).}
        \resumeItem{\tools{Rust, Python, Pest, SQLite, Streamlit, GitHub}}
      \resumeItemListEnd
\resumeSubHeadingListEnd

%-----------PROJETS-----------
\section{PROJETS}
    \resumeSubHeadingListStart
      \resumeProjectHeading{\textbf{Moteur de recherche scientifique}}{}
          \resumeItemListStart
            \resumeItem{Recherche d'information sémantique (TF-IDF, Sentence-BERT, LDA, graphe de citations) : F-mesure 0,847 / AUC 0,974 avec MPNet, vs 0,283 / 0,820 pour TF-IDF. }
            \resumeItem{\tools{Python, Sentence-BERT, scikit-learn, NetworkX}}
          \resumeItemListEnd
      \resumeProjectHeading{\textbf{Pipeline Reinforcement Learning \& MLOps}}{}
          \resumeItemListStart
            \resumeItem{REINFORCE (from scratch) et agents A2C (Stable-Baselines3) sur CartPole et bras robotique PandaReach : 100\% de succès ; modèles publiés sur Hugging Face, runs suivis sur W\&B. }
            \resumeItem{\tools{PyTorch, Stable-Baselines3, Gymnasium, W\&B}}
          \resumeItemListEnd
      \resumeProjectHeading{\textbf{Modèles génératifs : GANs \& Diffusion}}{}
          \resumeItemListStart
            \resumeItem{DCGAN/cGAN conditionnels (Pix2Pix), diffusion DDPM, et inférence Stable Diffusion 3.5 en quantization 4-bit. }\resumeItem{\tools{PyTorch, diffusers, bitsandbytes}}
          \resumeItemListEnd
    \resumeSubHeadingListEnd

%-----------FORMATION-----------
\section{FORMATION}
  \resumeSubHeadingListStart
    \resumeSubheading
      {École Centrale de Lyon (ECL)}{2024 -- 2026}
      {Diplôme d'ingénieur généraliste (double diplôme ENSI) }{Lyon, France}
    \resumeSubheading
      {École Nationale des Sciences Informatiques (ENSI)}{2022 -- 2024}
      {Ingénierie informatique --- Majore de promotion (200 étudiants), moyenne 16,01/20}{La Manouba, Tunisie}
    \resumeSubheading
      {Institut Préparatoire aux Études Scientifiques et Techniques (IPEST)}{2020 -- 2022}
      {Cycle préparatoire scientifique --- classée 66\textsuperscript{e} / 1500 au concours national}{La Marsa, Tunisie}
  \resumeSubHeadingListEnd

%-----------COMPETENCES & ACTIVITES-----------
\section{COMPÉTENCES \& CENTRES D'INTÉRÊT}
\small
\begin{minipage}[t]{0.60\textwidth}
     \textbf{Langages} : Python, C++, Java, SQL, NoSQL, JavaScript, Rust \\ \vspace{2pt}\textbf{APIs \& Frameworks} : REST, WebSocket, FastAPI, Flask, Node.js \\ \vspace{2pt}\textbf{Big Data \& Cloud} : Hadoop, Apache Spark, Azure, Docker,\\ Kubernetes, GitHub Actions \\ \vspace{2pt}\textbf{IA / Data} : PyTorch, keras, scikit-learn, pandas, Sentence-BERT, RL, \\ NLP, OR-Tools, Agentic AI, RAG, LLM \\ \vspace{2pt}\textbf{Data \& Infra} : MySQL, PostgreSQL, MongoDB, SQLite, Linux, Bash 
\end{minipage}\hfill
\begin{minipage}[t]{0.37\textwidth}
     \textbf{Clubs} : Responsable médias et design de l'ENSI Competitive Programming Club et du Forum ENSI édition 2023 (+5 passages radio/TV) \\ \vspace{2pt}\textbf{Sport} : Padel, randonnée \\ \vspace{2pt}\textbf{Loisirs} : Photographie, chant, lecture, films \& séries
     \\ \vspace{2pt}\textbf{Langues} : Arabe (natif), Français (bilingue), Anglais (C1), Espagnol (A2)
\end{minipage}

\end{document}
